% results_method

\section{Method}
\label{sec:method}

\subsection{Corpus and Data Collection}
\label{sec:method-corpus}

We compiled a corpus of user reviews for 26 Islamic prayer-time and devotional
applications on the Google Play Store, collected on 2026-08-19 and spanning
each application's review history available at that date. Reviews and app-level
metadata (store description, install count, rating) were scraped
programmatically. The resulting corpus comprises 732,194 reviews.
Figure~\ref{fig:corpus} summarises its composition.

Applications were selected purposively rather than by exhaustive enumeration of
the category. We prioritised high install and review counts, so that per-app
volume would support app-level comparison, and preferred applications whose
primary purpose is salat --- prayer times, adhan, qibla, tracking --- over
general Islamic-lifestyle or educational applications, retaining a
salat-specific application even where it offered fewer features overall. A
small number of lower-volume applications were kept where they were the only
instance of a functional category. The resulting set is neither a census of the
category nor a random sample, and no target size was fixed in advance; the
ecosystem-level claims in this paper are claims about these 26 applications,
which we take to cover the high-traffic core of the category rather than its
long tail.

One scoping decision deserves stating rather than leaving in the collection
script. Reviews were requested with the store locale set to English and the
United States (\texttt{lang=en}, \texttt{country=us}), sorted newest-first. This
does not yield a purely English corpus --- 48,265 of the 324,687 reviews
carrying text were subsequently identified as non-English
(\S\ref{sec:method-language}) --- but it does mean the corpus is the
English/US-locale slice of each application's reviews rather than its global
review population. Because every application was collected under identical
parameters, the cross-application comparisons throughout this paper remain
like-for-like; what the decision constrains is generalisation, and we return to
it in \S\ref{sec:limitations}.

\begin{figure}[t]
  \centering
  \includegraphics[width=\linewidth]{fig01_corpus_composition}
  \caption{Corpus composition: review volume for the 15 largest applications, language mix, and review-length distribution.}
  \label{fig:corpus}
\end{figure}

\subsection{Language Filtering and Analysis Sub-corpora}
\label{sec:method-language}

Reviews were filtered to those carrying review text, yielding an analysed
sub-corpus of 324,687 reviews (\(C_{\mathrm{text}}\)). Language identification
was then applied using fastText; the English-language subset used for topic
modelling comprises 276,422 reviews (\(C_{\mathrm{en}}\)).

\subsection{Sentiment Classification}
\label{sec:method-sentiment}

Reviews were scored in two passes: a lexicon-based pass with VADER and a
transformer-based pass with a RoBERTa sentiment classifier. The RoBERTa label is
the one reported throughout the paper, and VADER is retained as a comparison
point rather than as a source of any reported result.

Two separate checks bear on that choice, and we keep them apart because they
establish different things. First, agreement between methods across all 324,687
texted reviews (Figure~\ref{fig:sentiment-agreement}): RoBERTa agrees with the
star-derived label at Cohen's \(\kappa = 0.466\), against 0.210 for VADER, and
the two classifiers agree with each other at 0.385. Star ratings are not ground
truth --- RQ2 is in part an argument that they are not
(\S\ref{sec:rq2}) --- but RoBERTa tracking them more than twice as closely as
VADER does is the empirical reason we report the transformer label.

Second, we hand-labelled a 500-document gold set into four categories
(Positive, Neutral, Negative, Mixed); on the 100 rows coded by both annotators,
raw agreement was 89.0\% and Krippendorff's \(\alpha = 0.824\). We state plainly
what this does and does not establish: it is an inter-annotator reliability
figure showing that people apply the scheme consistently, not a measurement of
classifier accuracy against those labels, which we did not compute. Sentiment
therefore enters this paper as a comparatively, rather than absolutely,
validated variable, and the claims resting most heavily on it --- the
bimodality of tracker sentiment (\S\ref{sec:rq1}) and the feature-level
negative-sentiment rates behind RQ6a --- should be read with that gap in mind.

\begin{figure}[t]
  \centering
  \includegraphics[width=\linewidth]{fig02_sentiment_agreement}
  \caption{Pairwise agreement between the three sentiment signals across the 324,687 texted reviews, row-normalised, with Cohen's \(\kappa\) per pair. RoBERTa tracks the star-derived label more closely (\(\kappa = 0.466\)) than VADER does (\(\kappa = 0.210\)).}
  \label{fig:sentiment-agreement}
\end{figure}

\subsection{Aspect-Based Sentiment Lexicon}
\label{sec:method-aspects}

We built a lexicon of 27 aspects covering the functional and experiential
surface of devotional apps (prayer-time accuracy, adhan reminders, Quran audio,
qibla direction, prayer tracking, advertising, monetisation, and others),
comprising 340 keywords and 106 regular-expression patterns applied at the
sentence level. Each sentence can carry zero, one, or multiple aspect tags.
Keyword and pattern counts vary widely across the 27 aspects, and since several
results below turn on how thinly a given aspect is defined, we release the
lexicon in full rather than summarising it: every aspect's exact keyword list
and patterns are in the artefact (\S\ref{sec:method-repro}), as are the aspect
groupings each research question composes from them.

\subsection{Affect Lexicon}
\label{sec:method-affect}

RQ1 additionally uses two small keyword lists, one guilt- and anxiety-related
and one motivation-related, to compare affective register inside and outside
tracker discourse. These differ from every other measurement in this paper in
one respect we state rather than leave to be found: they carry no gold sample
and no precision figure, because the 300-row aspect draw and the 500-document
sentiment set label aspects and valence, neither of which is the guilt/motivation
distinction. We release both lists with the analysis code, report the result
they produce as a lexical-frequency comparison rather than as a validated
measurement of guilt (\S\ref{sec:rq1}), and rest the construct claim on the
hand-coded passes instead.

\subsection{Devotional-Formula Filtering}
\label{sec:method-devotional-filtering}

After lexicon tagging, we removed aspect assignments whose triggering sentence
consists of devotional formulae rather than substantive feedback: 1,686 of the
192,874 aspect-tagged sentence pairs the lexicon produced (0.9\%), leaving
191,188. The unit is the (aspect, sentence) pair rather than the sentence, so a
sentence carrying two tags is counted twice. This step is necessary because general-purpose English feature
vocabulary misfires badly on devotional register. The clearest case is
\texttt{goal\_system}, which scored 0\% precision against its gold rows because
reviewers routinely close a review with phrases such as ``May Allah reward
you,'' which the lexicon's goal-related keywords match as if they described a
tracking feature. We treat this as a substantive finding rather than a cleaning
footnote: off-the-shelf sentiment and aspect tooling, built and validated on
general product review text, does not transfer cleanly to religious register,
and a devotional-formula filtering pass is necessary before aspect-level
analysis is trustworthy in this domain.

\subsection{Demand Mining}
\label{sec:method-demand}

We extracted five types of demand signal from review text via pattern
matching, each attached to the aspect(s) it co-occurs with: \emph{request}
(explicit asks for a feature), \emph{absence} (noting a feature is missing),
\emph{churn} (citing a missing feature as a reason for leaving or switching
apps), \emph{preference} (citing a feature as the reason for choosing an app),
and \emph{tenure} (citing a feature in the context of long-term use). Against a
200-row gold sample stratified 40 rows per type, overall precision was 91.5\%
(183/200): preference 100\%, tenure 100\%, churn 97.5\%, request 95\%, and
absence 65\%.

This figure carries a provenance caveat that we state plainly rather than
suppress. The gold sample was drawn and labelled from one run of the demand
extractor; an upstream threshold subsequently changed, the extractor was
re-run, and the regenerated output no longer matches the labelled sample
row-for-row (only 11 of the 200 labelled reviewIds recur in the current
extraction). The demand counts reported elsewhere in this paper come from the
current extraction, while the 91.5\% precision figure was measured on its
immediate predecessor. We report it as the best available estimate of demand-pattern
precision, with this provenance gap disclosed, rather than as a measurement of
the exact rows underlying every reported count.

\subsection{Promise Extraction}
\label{sec:method-promise}

To compare what an application promises against what it delivers, we extracted
claimed features from each app's Google Play store description and compared
them against a hand-built feature matrix (one row per app, one column per
aspect). Agreement between the two sources is weak: mean Cohen's
\(\kappa = 0.272\) across 20 features, with a maximum of 0.62 for women's
tracking and four features at \(\kappa = 0.00\) --- travel concession,
calendar integration, prayer times by location, and timely reminders. We
report this plainly rather than treat the two sources as interchangeable;
Figure~\ref{fig:promise-agreement} gives the per-feature breakdown, including
cases where the description misses a feature the matrix records and cases where
it overclaims one the matrix does not.

Those four zeros need reading carefully rather than at face value. In each
case one source records the feature in almost every application or in almost
none, so there is too little variance for \(\kappa\) to correct against, and
the coefficient collapses to zero despite high raw agreement: 96.2\% for timely
reminders (26 of 26 against 25 of 26) and 96.2\% for travel concession. Mean
raw agreement across the 20 features is 73.1\%. We report both summaries
because \(\kappa\) alone would understate how far the two sources actually
coincide on the features RQ6a and RQ3 depend on, while raw agreement alone
would flatter the features where the sources genuinely diverge.

\begin{figure}[t]
  \centering
  \includegraphics[width=\linewidth]{fig03_promise_source_agreement}
  \caption{Agreement between store-description promises and the hand-built feature matrix, by feature.}
  \label{fig:promise-agreement}
\end{figure}

\subsection{Topic Modelling}
\label{sec:method-topics}

We fit a BERTopic model over \(C_{\mathrm{en}}\) (276,422 documents) using a
seeded UMAP reduction and a stopword list extended with Islamic-domain terms
and application-name tokens, to prevent generic devotional vocabulary and
app-name mentions from dominating topic keywords. The ecosystem-level model
produced 40 topics plus an outlier bucket. Topic 0 is an undifferentiated
generic-praise cluster (168,054 documents, 61\% of modelled documents) and the
outlier bucket holds a further 62,745 (22.7\%); the remaining 39 topics, which
we label and report, cover approximately 45,600 documents (about 16\% of
\(C_{\mathrm{en}}\)). All 63 topics across the ecosystem model and three
targeted sub-models (described below) were labelled by hand across four
sheets, alongside each model's outlier bucket. We treat the topic model as characterising the interpretable minority
of the corpus, not the corpus as a whole, throughout the paper.

In addition to the ecosystem-level model, we fit three targeted sub-models on
aspect-restricted subsets to support RQ1 (tracker discourse), RQ4 (menstrual
and inclusion discourse), and RQ5 (feature-bloat discourse). Topic structure
under BERTopic is sensitive to the UMAP seed; we assessed this by refitting
each sub-model across five seeds fixed in advance (42--46), before inspecting
any result. The \texttt{rq1\_tracker} sub-model is stable (mean pairwise
Adjusted Rand Index 0.837); \texttt{rq5\_bloat} is not (6--11 topics across the
five seeds, mean pairwise ARI 0.452). Where we draw on \texttt{rq5\_bloat}, we
therefore claim the recurring themes rather than a specific partition into a
fixed number of topics.

\subsection{Statistical Models}
\label{sec:method-models}

We fit five models (M1--M5), one primary per research question for RQ1--RQ5:
M1 (RQ2, the competing costs of accuracy and interface complaints), M2 (RQ1,
sentiment by tracker tier), M3 (RQ4, inclusion), M4 (RQ5, feature bloat and
hardware), and M5 (RQ3, competitive edge features). RQ6 is answered
descriptively, by comparing promised features against the sentiment their
mentions record, and enters no inferential model. The primary specification for
each of the five is a review-level mixed-effects
model with a random intercept per application, which accounts for the fact
that reviews cluster within 26 apps of very different size and character. Where
a mixed model did not converge, we report a cluster-robust OLS fit on the same
data and specification as a sensitivity estimate, clustering standard errors on
application.

Coefficients from the primary specifications enter a single Benjamini--Hochberg
false-discovery-rate correction: 29 in total, of which 20 survive at
\(q < 0.05\). The family's composition needs stating precisely, because it
treats two kinds of refit differently. Cluster-robust refits are reported
alongside their primary counterparts and are not entered, since correcting a
robustness check against the test it checks would double-count the contrast.
The eight ordinal refits of M1's complaint terms \emph{are} entered, because
they change the outcome variable rather than only the error structure and so
constitute separate tests. We flag the asymmetry rather than leave it to be
inferred: it is a defensible line but not the only defensible one, and a reader
who would draw it elsewhere can recompute from the released coefficient
table.

Eleven terms cannot support inference, and we report all eleven rather than
silently dropping them; they divide into two kinds that sit differently with
respect to that denominator. Eight are mixed-model terms whose standard errors
are degenerate rather than merely large --- M2's two tracker-tier terms and six
from M4, with standard errors from \(4.1\times10^{2}\) to \(2.0\times10^{6}\)
and every one reporting \(p \approx 0.99999\). These eight do carry
\(q\)-values and are inside the 29, so read in isolation they resemble
``tested, not significant'' rather than ``failed to converge''; they account
for eight of the nine terms in the 29 that fail to survive correction, the
ninth being \texttt{madhab} in the ordinal specification, which we enter as a
primary term rather than a sensitivity refit and which returns a small,
imprecise estimate that reverses sign against M1 (\S\ref{sec:rq2}); we rest no
directional claim on it. The other three ---
M3's \texttt{women\_feature\_score\_diff}, M5's \texttt{has\_mosque\_finder},
and one M5 demand term --- return neither a standard error nor a \(p\)-value,
so they are outside the 29 rather than inside it, and no correction applies to
them at all. We
attribute the non-convergence to a structural cause: a random intercept per app,
fitted over only 26 apps, cannot be identified separately from an app-level
predictor in the same model. Throughout, we cite null results only from a model
that converged, and never describe one of these eleven as evidence of no effect.

\subsection{Measurement Quality}
\label{sec:method-measurement}

Because several of the paper's claims rest on aspect tags, we report their
precision directly rather than assuming lexicon output equals ground truth.
Against a 300-row gold sample, of which 192 rows still match under the current
lexicon configuration, unweighted aspect precision is 69.3\% and
corpus-volume-weighted precision is 84.4\% (95\% CI [79.7, 88.9]). The gold set
was originally labelled against an earlier, broader configuration including a
zero-shot backstop; every subsequent lexicon change has only removed matches,
never added them, so the 192 surviving rows are an unbiased subset of the
original 300 rather than a re-labelling exercise. Two aspects,
\texttt{calendar\_sync} and \texttt{goal\_system}, have no surviving gold rows
(together 0.2\% of aspect-tagged volume) and are excluded from the weighted
figure.

The weighted figure is driven by high-volume aspects and conceals substantial
variation. The 300-row gold draw sampled all 27 aspects at 11 rows each.
Table~\ref{tab:aspect-precision} reports precision for the 25 whose rows
survive the current lexicon configuration; \texttt{calendar\_sync} and
\texttt{goal\_system} lost all of theirs to lexicon narrowing. The 25 account
for all 192 surviving rows and for the weighted figure above. Seven of them
fall below 50\%, with gold samples as small as one or two rows. We report this
in full because several of these are load-bearing for individual research
questions: every aspect-level count drawn from one of the seven is reported
elsewhere as a precision-adjusted estimate with its interval, never as a bare
tagger count. Seventeen aspects sit at or above 64\% precision, and these
supply almost all model input, with two exceptions we flag rather than bury.
\texttt{madhab}, at 60\%, enters M1, and its coefficient is also the one M1
term that reverses sign under the ordinal robustness check, so we rest no
directional claim on it (\S\ref{sec:rq2}). And the mosque-finder contrast in
RQ3 is fit on the subset of reviews the \texttt{mosque\_finder} tagger selects,
which runs at 27\%; because the resulting noise falls on both arms of that
comparison rather than one, we read the direction of the gap and not its
magnitude, and say so at the point of use (\S\ref{sec:rq3}). No other model
term is drawn from an aspect below 50\% precision.

\begin{table*}[t]
  \centering
  \small
  \caption{Aspect-tagging precision against gold-standard rows, by aspect, for the 25 aspects whose gold rows survive the current lexicon configuration. Bold rows fall below 50\% precision; any count for these is reported elsewhere only as a precision-adjusted estimate with its interval. Gold samples are small, so the intervals are wide and we report them rather than the point estimate alone.}
  \label{tab:aspect-precision}
  \begin{tabular}{lrrrll}
    \toprule
    Aspect & Corpus \(n\) & Gold \(n\) & Precision & 95\% CI & Adjusted estimate \\
    \midrule
    reminders\_adhan        & 30{,}069 & 7  & 100\% & [65, 100] & \\
    quran\_audio            & 24{,}964 & 6  & 100\% & [61, 100] & \\
    prayer\_times\_accuracy & 21{,}941 & 11 & 73\%  & [43, 90]  & \\
    ads\_intrusive          & 17{,}845 & 11 & 100\% & [74, 100] & \\
    ui\_design              & 15{,}720 & 11 & 73\%  & [43, 90]  & \\
    monetization            & 13{,}914 & 11 & 82\%  & [52, 95]  & \\
    adhkar                  & 9{,}280  & 3  & 100\% & [44, 100] & \\
    qibla                   & 8{,}526  & 11 & 82\%  & [52, 95]  & \\
    stability\_bugs         & 7{,}665  & 11 & 64\%  & [35, 85]  & \\
    privacy\_data           & 7{,}471  & 10 & 90\%  & [60, 98]  & \\
    prayer\_tracker         & 6{,}009  & 9  & 78\%  & [45, 94]  & \\
    complexity\_bloat       & 5{,}389  & 9  & 67\%  & [35, 88]  & \\
    widgets                 & 4{,}755  & 11 & 100\% & [74, 100] & \\
    spiritual\_affect       & 3{,}395  & 2  & 100\% & [34, 100] & \\
    \textbf{mosque\_finder}      & 3{,}043 & 11 & \textbf{27\%} & [10, 57] & \(\approx\)830 [297, 1{,}721] \\
    \textbf{companion\_hardware} & 2{,}512 & 5  & \textbf{40\%} & [12, 77] & \(\approx\)1{,}005 [295, 1{,}932] \\
    \textbf{women\_period}       & 2{,}194 & 10 & \textbf{20\%} & [6, 51]  & \(\approx\)439 [124, 1{,}119] \\
    \textbf{forbidden\_times}    & 1{,}534 & 10 & \textbf{10\%} & [2, 40]  & \(\approx\)153 [27, 620] \\
    \textbf{table\_format}       & 1{,}483 & 9  & \textbf{44\%} & [19, 73] & \(\approx\)659 [280, 1{,}088] \\
    nafl\_times             & 963      & 5  & 100\% & [57, 100] & \\
    calc\_method            & 639      & 10 & 90\%  & [60, 98]  & \\
    \textbf{qasr\_travel}        & 637      & 2  & \textbf{0\%}  & [0, 66]  & do not report volume \\
    madhab                  & 491      & 5  & 60\%  & [23, 88]  & \\
    guides                  & 202      & 1  & 100\% & [21, 100] & \\
    \textbf{tracker\_score}      & 101      & 1  & \textbf{0\%}  & [0, 79]  & 1 gold row, uninformative \\
    \midrule
    Total                   & 190{,}742 & 192 & & & \\
    \bottomrule
  \end{tabular}
\end{table*}

Gold-sample size per aspect ranges from one to eleven rows, so these intervals
are wide; we treat the width as honest uncertainty rather than a reason to
report the point estimate alone. We also measured a zero-shot classification
backstop intended to supplement the lexicon, and dropped it after it scored
11.5\% precision on 300 hand-labelled sentences; we report this discarded
component as evidence of calibration during development rather than omit it.

\subsection{Research Ethics and Data Handling}
\label{sec:method-ethics}

The corpus is public store review text. No user was recruited, contacted, or
observed, and no account-linked or private data was accessed, so the study
involved no human participants and was not submitted for human-subjects
review. We none the less treat it as
data about people: the reviews concern religious observance, and some disclose
menstruation, illness, or lapses in practice, which a reviewer complaining about
a notification may not have understood themselves to be publishing. We
therefore report no reviewer names or identifiers, attribute quotations only to
an application and a star rating, and choose quotations carrying no incidental
personal detail. We quote verbatim rather than paraphrase, since altering
wording in a study whose claim concerns devotional register would misrepresent
the evidence, and we accept the residual risk that a short quoted string can be
matched back through store search. Two of our quotations disclose menstruation, the category we name above as
most sensitive. We judged the disclosure necessary: RQ4's claim is that users
voice this need explicitly, and paraphrase would convert direct evidence into
our own assertion of it. We chose the shortest instances that carry the point,
neither names an author, and both are already public and searchable. The
artefact releases derived data rather than redistributable raw review text
(\S\ref{sec:method-repro}). We name
applications throughout, including in criticism: they are public commercial
products, and the argument cannot be made without saying which application did
what. Where we name one in connection with a specific controversy rather than
an ecosystem pattern, we confine the claim to that episode
(\S\ref{sec:limitations}).

\subsection{Use of AI Assistance}
\label{sec:method-ai}

All analysis software for this study was written with the assistance of Claude
Code, Anthropic's command-line coding agent, driven interactively by the
authors. We disclose this because the assistance was substantive in the conduct
of the research rather than confined to preparing the manuscript. Three
boundaries bound what it can threaten. First, no generative model produced,
imputed, translated, or altered any review text, any label used as data, or any
reported statistic; every number here is the output of code executed over the
scraped corpus. Second, the design is the authors' own: the six research
questions were specified in a written analysis plan committed to version
control before the analysis was run, the aspect lexicon and the coding frame
were hand-built, and every gold-standard label was assigned by a human
annotator without model assistance. Third, the authors read, executed, and are
responsible for all code.

Assistant-written code can carry defects a human author would not introduce,
and we do not claim otherwise; our safeguards are the ones already reported ---
per-stage precision measured against hand-labelled gold sets, a component those
checks rejected at 11.5\% precision, the seed-stability pass, and the
byte-identical re-run below. One such defect was found and fixed during
analysis: the three sub-models shared a hardcoded output path and overwrote the
ecosystem topic sheet on each run. The sheet was rebuilt from the saved model
and stored assignments without refitting, and no reported statistic depended on
it. Confidence in the pipeline should rest on the checks that catch such
faults, not on an assumption that none occurred.

\subsection{Reproducibility}
\label{sec:method-repro}

The analysis reported in this paper was run on 2026-09-03 at commit
\texttt{5da2db7}. We subsequently re-ran the modelling stage from a clean tree
at the same commit and obtained byte-identical output, confirming
reproducibility of the reported coefficients. All stochastic components are
seeded: the UMAP reduction underlying topic modelling is seeded, and the five
seeds used for the topic-stability check (42--46) were fixed in advance of
inspecting any result, so the seed set could not be adjusted after the fact to
favour a particular outcome. A run manifest recording commit, hardware, batch
size, and library versions for every reported figure accompanies the research
artefact.

\paragraph{Data and code availability.} We release the full pipeline, the
aspect lexicon, the feature matrix, every hand-labelled gold and topic sheet,
the model output, and the run manifest, at a URL supplied on acceptance and
withheld here for anonymous review. Derived data sufficient to recompute every
table and figure is included. Two things are not: the raw review corpus, which
the store's terms do not permit us to redistribute and which carries the
re-identification risk noted in \S\ref{sec:method-ethics}, and the fitted model
binaries, which are large and reproducible from the seeded pipeline. The gold
sheets we do release carry the review text their labels apply to, since a
precision figure cannot be checked without it; that is a few hundred rows drawn
for verification rather than the corpus, and we judge the disclosure
proportionate on the same reasoning that governs quotation above. The
collection script and application list are included, so the corpus can be
re-gathered, though store histories change and an exact reconstruction of our
snapshot is not possible.
